\WorkEntry{ATH-WORK-0008}{Gyre}{%
\begin{OldSchoolEssay}
Put a blind man in a field and ask him to walk straight ahead, and yes, he walks a circle. They told me with authority that this raw instinct brings him back to his starting point. Some primordial instinct---a self-preservation behaviour born out of evolutionary advantage to take him back to his cave; he's a homing pigeon. We should all build a cave dwelling for our Alzheimer's sufferers, right?. We bought into that stuff back in the 70's---an 'authority' told us so. Did you try it? I did. My circle was anticlockwise, just as most of our testicles hang lower on the left, and our breasts swell more on that side too. Correlation is there in the old studies. The fault here is the assumption that without excitement (stimulus) otherwise, we walk a straight path---we don't; we walk in circles.

The second law of thermodynamics states that order is unnatural and everything tends to disorder. What? The universe, galaxies, solar systems, and so forth are order---entropy defiled. A paradox?

We're stuck in cycles. Gravity binds. Love binds. Disorder is abhorrent. We build and build; our lives, our loves, our families, our homes---we take chaos and order it. That second law? We're always fighting it. The habitual smoker, the habitual drinker, the habitual worker---all of us---we are bound to cycles. But remember, quick victories lead to longer-term defeats.

Love is order. But you see, the building of structure from chaos does not actually fly in the face of entropy as we might think---in fullness, we serve entropy with an overflowing plate. Remember, building relationships (or structure) takes work, and work is energy. Energy dissipated is entropy served. All structures crumble, eventually returning to chaos. Entropy is patient.

We circle back to what we know, again and again, patterns and forms universally repeated. Chaos is the fabric of our universe---bits and pieces plopping into and out of existence---yet somehow it supports our ethereal existence.

The manifestation of the quantum, our reality, has to marshal energy carefully, and with anything but randomness. Efficiency is paramount, and supremely served by this repetitious use of the same patterns and forms. The structure of our materials, temporary as they are, must be efficiently built and supported. What's the point of a physical manifest without lifetime enough to support a purpose?

Circles.

To keep doing the same thing, from the micro to the macro, is the most efficient thing to do. A perfect circle; a perfect sphere; those are the structures that, out of all, are the most energy-efficient and therefore (almost by design) the longest living. Decay is rebuffed for as long as it can be.

This is why we live our lives in circles.
\end{OldSchoolEssay}
}{\VLabel{Lens:}  Repetition is efficiency.}{\VLabel{Link:}  Habit / Efficiency / Entropy.}{\VLabel{Question:}  Which circle are you mistaking for destiny?}
